% ABSTRACT

{
	\cleardoublepage
	\let\clearpage\relax
	\noindent SOUSA, M. A. B., \textbf{Implementation of a dynamic backlash model in gears for rotational dynamics analysis} 2026. 92 p. Master Dissertation, Federal University of Uberlândia, Uberlândia.
	
	\chapter*{\MakeUppercase{Abstract}}
}

%Gearboxes are essential components in industrial rotating machinery, but non-linear phenomena such as dynamic \textit{backlash} and Time-Varying Mesh Stiffness (TVMS) introduce high complexity to their global structural response. Thus, it becomes necessary to develop tools capable of adequately representing the system's behavior and assisting in the understanding of the phenomena involved. In this context, the objective of this study was to develop a mathematical model capable of interpreting such effects and relating them to rotordynamics. To this end, the non-linear gear mesh model by \citeonline{yi2019} was coded in Python and integrated into the \textit{Rotordynamic Open-Source Software} (ROSS) library, based on the Finite Element Method (FEM). The original formulation of 3 degrees of freedom (\textit{DoF}) per gear was expanded to 6 \textit{DoF} per node, enabling the inclusion of geometric parameters and helical gear forces, based on \citeonline{kubur2004} and \citeonline{mo2025}. Additionally, a smoothing technique using hyperbolic tangents \cite{walha2006} was incorporated to handle numerical discontinuities and optimize the convergence of time integrators. To solve the differential equations of motion, the Newmark method was used in combination with the Newton-Raphson algorithm and an adaptive time step. A simplified analytical model was also developed to determine the natural frequencies of the geared pair. The results were validated against the literature, showing natural frequencies perpendicular and parallel to the Line of Action (LOA). Bifurcation diagrams confirmed clear transitions between periodic and chaotic motion regimes at high speeds. Finally, the integrated methodology was applied to a real practical case of a large-scale industrial gearbox operated by Petrobras (input speed of 1953 rpm and output speed of 12225 rpm) employing the complete gear mesh stiffness profile calculated by ROSS. For this case, the smoothing tool indicated a reduction of approximately $7.3\%$ in the total computational processing time without loss of the dynamic system's physical representation. It is concluded that the developed approach constitutes a robust and high-fidelity solution for mechanical design and fault diagnosis in complex industrial transmissions.

Gearboxes are essential components in industrial rotating machinery, but nonlinear phenomena such as dynamic backlash and time-varying mesh stiffness (TVMS) introduce high complexity to their overall structural response. Thus, it becomes necessary to develop tools capable of adequately representing system behavior and aiding in the understanding of the involved phenomena. In this context, the objective of this study was to develop a mathematical model capable of interpreting such effects and relating them to rotational dynamics. To this end, a nonlinear meshing model was coded in Python and integrated into the \textit{Rotordynamic Open-Source Software} (ROSS) library, based on the Finite Element Method (FEM). The original formulation of 3 degrees of freedom (DOF) per gear was expanded to 6 DOF per node, enabling the inclusion of geometric parameters and helical gear forces. Additionally, a hyperbolic tangent smoothing technique was incorporated to handle numerical discontinuities and optimize the convergence of time integrators. To solve the differential equations of motion, the Newmark method combined with the Newton-Raphson algorithm and adaptive time stepping was used. A simplified analytical model was also developed to determine the natural frequencies of the gear pair. The results were validated against literature, showing natural frequencies perpendicular and parallel to the Line of Action (LOA). Bifurcation diagrams confirmed clear transitions between periodic and chaotic motion regimes at high speeds. Finally, the integrated methodology was applied to a real practical case of a large-scale industrial gearbox operated by Petrobras (input rotation of 1953~rpm and output of 12225~rpm) employing the full mesh stiffness profile calculated by ROSS. For this case, the smoothing tool indicated a reduction of approximately $7.3\%$ in total computational processing time without loss of physical representativeness of the dynamic system. It is concluded that the developed approach represents a solution for mechanical design and fault diagnosis in complex industrial transmissions.

\vspace*{\fill}
\noindent\rule{\textwidth}{1pt}
\textit{Keywords: Backlash, Gear Mesh Stiffness, Gear, Dynamic Analysis, ROSS}